% =====================================================================
% propositions.tex -- two formal results for paper.tex
%
% PREAMBLE LINES NEEDED (paper.tex loads amsmath, amssymb only):
%   \usepackage{amsthm}
%   \newtheorem{proposition}{Proposition}
%   \newtheorem{lemma}{Lemma}
%   \theoremstyle{remark}
%   \newtheorem{remark}{Remark}
%   \newcommand{\Var}{\operatorname{Var}}
%   \newcommand{\vc}{\textsuperscript{\,[verify citation]}}   % set to {} once checked
% The paper already defines \E, \Cov, \HSIC, \MMD, \one.
% New bib keys are in theory/refs-additions.bib. Every one is marked \vc.
%
% Suggested placement (see theory/NOTES.md):
%   Statements: new Section 2.4 "Two exact results", after "Independence
%   statistics" and before Table 1 (tab:scope). Add two rows to Table 1.
%   Proofs and lemmas: new appendix "Proofs of Propositions 1 and 2".
%   Cross-references: Section 4.3 (first finding, SEAS-M) and Section 4.4
%   (third pattern, split-MMD mean shift).
% =====================================================================

% ---------------------------------------------------------------------
% PART 1. STATEMENTS (for Section 2.4)
% ---------------------------------------------------------------------

\subsection{Two exact results}\label{sec:props}
Two of the empirical findings follow from exact calculations. The first explains why a seasonal mean with a random phase misleads single-path tests. The second explains why memory removes the power of split MMD against a mean shift.

Throughout, $\xi$ and $\zeta$ are independent unit-variance fGn sequences with Hurst parameter $H\in(0,1)$. A \emph{single-path test} is a test whose decision is a measurable function of one observed pair $(X_t,Y_t)_{t=1}^n$, possibly with auxiliary randomness such as multiplier draws. All its reference values are computed from that same pair. The size of such a test is its rejection probability under the joint law of the phases and the noise. This is the quantity that EXP2 estimates.

\paragraph{Seasonal mean with a random phase.}
Fix an amplitude $a>0$ and an integer period $P\ge3$. Let $U$ and $V$ be independent and uniform on $\{0,\ldots,P-1\}$, and independent of $(\xi,\zeta)$. Set
\begin{equation}\label{eq:seasm}
X_t=\xi_t+a\sin\{2\pi(t+U)/P\},\qquad Y_t=\zeta_t+a\sin\{2\pi(t+V)/P\}.
\end{equation}
The processes $X$ and $Y$ are independent. EXP2 uses $a=1$ and $P=50$. Let $D=(U-V)\bmod P$ be the phase gap. It is uniform on $\{0,\ldots,P-1\}$. Define
\begin{equation}\label{eq:climit}
c(d)=\frac{(a^2/2)\cos(2\pi d/P)}{1+a^2/2}.
\end{equation}
Let $r_n$ be the sample Pearson correlation and $\hat C_n=n^{-1}\sum_t(X_t-\bar X)(Y_t-\bar Y)$ the sample covariance.

\begin{proposition}[Seasonal mean]\label{prop:seasm}
Under \eqref{eq:seasm}, the following hold.
\begin{enumerate}[label=(\roman*),nosep]
\item $r_n\to c(D)$ almost surely. The limit is random and has mean zero. Moreover $\Pr\{c(D)=0\}=2/P$ if $4$ divides $P$, and $0$ otherwise.
\item Consider any test that rejects when $|\hat C_n|>z_n(\hat v_n/n)^{1/2}$, where $z_n$ is bounded and $\hat v_n\ge0$ satisfies $\hat v_n/n\to0$ in probability. Its size satisfies $\liminf_n\Pr(\text{reject})\ge\Pr\{c(D)\neq0\}$, which equals $1$ for $P=50$. The Newey--West Bartlett rule with bandwidth $S_n=o_p(n)$ and the equal-weighted cosine rule with $\nu_n\to\infty$ projections belong to this class. So do the implementations used here.
\item Let $T_n=nV_n$ be the Gaussian-kernel HSIC statistic of Section~\ref{sec:methods}, with median-distance bandwidths, and let $P\ge5$. Then $V_n\to h(D)$ almost surely for a function $h$ with $\min_d h(d)>0$. Any test that rejects when $T_n>\kappa_n$, with $\kappa_n/n\to0$ in probability, therefore has size tending to $1$. The window-subsampling rule with window length $\ell_n=o(n)$ belongs to this class. The block-Markov wild rule with $\ell_n=o(n)$ and a fixed number $B\ge\alpha^{-1}-1$ of draws, centred or raw, also has size tending to $1$.
\end{enumerate}
\end{proposition}

The proposition does not say that every single-path test fails. It covers tests whose critical value is of smaller order than the statistic's seasonal component. For Pearson that means a critical value for $r_n$ that tends to zero. For HSIC it means a critical value for $T_n$ of order $o(n)$. The circular-shift rule is not in this class, because its reference values are themselves of order $n$ (Remark~\ref{rem:shift}).

The limit also explains the finite-sample rate. With $a=1$ and $P=50$, the gaps $d\in\{12,13,37,38\}$ give $|c(d)|=0.021$, and every other gap gives $|c(d)|\ge0.06$. At $n=2000$, $\sqrt n\,|c(d)|\approx0.9$ at those four gaps, so a $5\%$ test rarely rejects there. Those gaps carry probability $4/50$, which predicts a rejection rate near $92\%$. EXP2 records $91.4\%$ for Pearson NW. The HSIC limit $h(d)$ is also smallest at these gaps.

\paragraph{Split-sample mean shift.}
Let $n=2m$ and $Z_t=\mu+\xi_t+\delta\,\one\{t>m\}$, with $\mu$ unknown. Write $\Gamma$ for the $n\times n$ fGn covariance matrix, $\one$ for the vector of ones, and $c=(0,\ldots,0,1,\ldots,1)^\top$ for the second-half indicator. Let $z_q$ be the standard normal $q$-quantile and $\Phi$ its distribution function.

\begin{proposition}[Half-mean spread]\label{prop:halfmean}
\begin{enumerate}[label=(\roman*),nosep]
\item Let $\Delta_m=m^{-1}\sum_{t\le m}\xi_t-m^{-1}\sum_{t>m}\xi_t$. Then, exactly for every $m$,
\begin{equation}\label{eq:halfvar}
\Var(\Delta_m)=\frac{4m^{2H}-(2m)^{2H}}{m^2}=(4-2^{2H})\,m^{2H-2}.
\end{equation}
Hence the half-mean contrast has signal-to-spread ratio $R_m=\delta\,m^{1-H}/(4-2^{2H})^{1/2}$.
\item The two-sided level-$\alpha$ test that rejects when the observed half-mean difference $|\bar Z_{1:m}-\bar Z_{m+1:2m}|=|\Delta_m-\delta|$ exceeds $z_{1-\alpha/2}\Var(\Delta_m)^{1/2}$ has power $\Phi(R_m-z_{1-\alpha/2})+\Phi(-R_m-z_{1-\alpha/2})$.
\item Let $\phi$ be any test, possibly randomized, that is invariant under $Z\mapsto Z+b\one$ for every real $b$, and has level $\alpha$ when $\delta=0$. Then its power at $\delta>0$ is at most $1-\Phi(z_{1-\alpha}-\delta I_n^{1/2})$, where
\begin{equation}\label{eq:info}
I_n=c^\top\Gamma^{-1}c-\frac{(\one^\top\Gamma^{-1}c)^2}{\one^\top\Gamma^{-1}\one}
\end{equation}
is the inverse variance of the generalized-least-squares estimator of $\delta$. The matrix $\Gamma$ is invertible because the fGn spectral density is positive.
\end{enumerate}
\end{proposition}

At $n=2000$ and $\delta=0.1$, part (i) gives $R_m=2.24$ at $H=0.5$, $0.68$ at $H=0.7$, and $0.40$ at $H=0.8$. Part (ii) then gives a power of $61\%$, $10.5\%$, and $7\%$ for the linear half-mean test at $5\%$. A direct Toeplitz solve gives $\delta I_n^{1/2}=2.24$, $0.71$, and $0.46$. Part (iii) therefore caps the power of every location-invariant $5\%$ test at $72\%$, $18\%$, and $12\%$. Gaussian-kernel split MMD is location-invariant. Its kernel and median bandwidth depend on the data only through differences $Z_s-Z_t$. The wild, subsampling, and self-normalized rules are built from the same kernel values, and the size-adjusted cutoff is a constant. Every MMD rule in Table~\ref{tab:power} is therefore bound by part (iii), up to Monte Carlo error in the size adjustment. The observed $6.4$--$7.5\%$ at $H=0.8$ lies under the $12\%$ cap, and the $7$--$9\%$ at $H=0.7$ lies under the $18\%$ cap. The loss of power is a loss of information in the data, not a failure of a kernel rule.

Part (iii) is a bound, not a description of the Gaussian kernel. A heuristic links the two. The rank-one Hermite coefficient of the feature map is $a_1=\E\{\xi_0\phi(\xi_0)\}$. By Stein's identity it equals the derivative of the mean embedding of $\xi_0+\delta$ at $\delta=0$. To first order in $\delta$, the shift therefore moves the embedding along $a_1$, and the rank-one noise along $a_1$ is $\Delta_m a_1$. Along this direction the signal-to-spread ratio is $R_m$. A Hermite term of rank $r$ adds noise of variance order $m^{\max\{r(2H-2),-1\}}$, up to a logarithm at equality, in other directions. At $H=0.5$ they have the same order as the rank-one term, which may explain why oracle MMD ($45.7\%$) falls short of the linear test ($61\%$). At $H=0.8$ the rank-one term dominates and the two agree ($6.4\%$ against $7\%$). We have not proved this expansion for the RKHS-valued statistic.

% ---------------------------------------------------------------------
% PART 2. PROOFS (for a new appendix)
% ---------------------------------------------------------------------

\section{Proofs of Propositions~\ref{prop:seasm} and~\ref{prop:halfmean}}\label{sec:proofs}

\subsection{Ergodic averages with periodic weights}
Write $W_t=(\xi_t,\zeta_t)$. Let $\gamma(k)=\tfrac12(|k+1|^{2H}-2|k|^{2H}+|k-1|^{2H})$ be the fGn autocovariance. For $H\ne1/2$, $\gamma(k)\sim H(2H-1)k^{2H-2}$, and for $H=1/2$, $\gamma(k)=0$ when $k\ne0$. In both cases $\gamma(k)\to0$ because $H<1$.

\begin{lemma}\label{lem:ergodic}
Fix $L\ge0$ and $P\ge1$. Let $f_0,\ldots,f_{P-1}$ be measurable functions of $(W_t,\ldots,W_{t-L})$ with $\E|f_j(W_L,\ldots,W_0)|<\infty$. Then, almost surely,
\[
\frac1n\sum_{t=L}^{n+L-1}f_{t\bmod P}(W_t,\ldots,W_{t-L})\longrightarrow\frac1P\sum_{j=0}^{P-1}\E f_j(W_L,\ldots,W_0).
\]
\end{lemma}

\begin{proof}
A stationary Gaussian sequence whose covariance tends to zero is mixing under the shift \citep{maruyama1949harmonic}\vc. See also \citet[Chapter~14]{cornfeld1982ergodic}\vc. So the shifts of $\xi$ and of $\zeta$ are mixing. The two sequences are independent, and the direct product of two mixing systems is mixing \citep[Section~1.7]{walters1982introduction}\vc. Hence the shift $S$ of $W$ is mixing. Then $S^P$ is mixing, because $\Pr(S^{-Pk}A\cap B)\to\Pr(A)\Pr(B)$ is a subsequence of the mixing limit for $S$. A mixing transformation is ergodic. For each residue $j$, Birkhoff's theorem applied to $S^P$ and to $f_j$ evaluated at the window ending at time $j+L$ gives $K^{-1}\sum_{k<K}f_j(W_{j+L+kP},\ldots,W_{j+kP})\to\E f_j$ almost surely. The sum in the lemma splits into these $P$ residue classes. Each class holds $n/P+O(1)$ terms. Averaging the $P$ limits with weights tending to $1/P$ gives the claim.
\end{proof}

Two special cases are used repeatedly. Taking $f_j(w)=g_j\,f(w)$ gives $n^{-1}\sum_t g_t f(W_t)\to\bar g\,\E f(W_0)$ for a $P$-periodic sequence $g$ with mean $\bar g$. Taking $f_j(w)=h\{w+(s_j,q_j)\}$ gives limits of empirical averages of a periodically shifted sequence.

\subsection{Proof of Proposition~\ref{prop:seasm}(i)}
The phases $(U,V)$ are independent of $W$ and take finitely many values. It is therefore enough to prove almost-sure convergence for each fixed pair $(u,v)$. Write $s_t=a\sin\{2\pi(t+u)/P\}$ and $q_t=a\sin\{2\pi(t+v)/P\}$. For $P\ge3$ the trigonometric sums give
\[
\frac1P\sum_{j=0}^{P-1}s_j=0,\qquad\frac1P\sum_{j=0}^{P-1}s_j^2=\frac{a^2}2,\qquad\frac1P\sum_{j=0}^{P-1}s_jq_j=\frac{a^2}2\cos\frac{2\pi(u-v)}P.
\]
The last identity uses $\sin\alpha\sin\beta=\tfrac12\{\cos(\alpha-\beta)-\cos(\alpha+\beta)\}$. The sum of $\cos(\alpha_j+\beta_j)$ vanishes because $e^{4\pi i/P}\ne1$ for $P\ge3$. Lemma~\ref{lem:ergodic} with $L=0$ gives, almost surely, $\bar X\to0$ and $\bar Y\to0$. It also gives
\[
\frac1n\sum_tX_t^2=\frac1n\sum_t\xi_t^2+\frac2n\sum_ts_t\xi_t+\frac1n\sum_ts_t^2\to1+0+\frac{a^2}2.
\]
The same holds for $Y$. For the cross moment, $n^{-1}\sum_t\xi_t\zeta_t\to\E\xi_0\zeta_0=0$, and the terms $n^{-1}\sum_ts_t\zeta_t$ and $n^{-1}\sum_tq_t\xi_t$ tend to $0$. Hence $\hat C_n\to(a^2/2)\cos\{2\pi(u-v)/P\}$ and $r_n\to c(u-v)$.

The gap $D=(U-V)\bmod P$ is uniform because $U$ is uniform and independent of $V$. So $\E\cos(2\pi D/P)=P^{-1}\sum_d\cos(2\pi d/P)=0$. Finally, for $0\le d<P$, $\cos(2\pi d/P)=0$ exactly when $d=P/4$ or $d=3P/4$. Such integers exist only when $4$ divides $P$. \qed

\subsection{Proof of Proposition~\ref{prop:seasm}(ii)}
Let $C(D)=(a^2/2)\cos(2\pi D/P)$, so $\hat C_n\to C(D)$ almost surely by part (i). The variable $D$ takes finitely many values. Choose $\eta>0$ below $\tfrac12\min\{|C(d)|:C(d)\ne0\}$. On $\{C(D)\ne0\}$, the test fails to reject only if $|\hat C_n|\le2\eta$ or $z_n(\hat v_n/n)^{1/2}\geq 2\eta$. The first event has probability tending to zero by almost-sure convergence. The second has probability tending to zero because $z_n$ is bounded and $\hat v_n/n\to0$ in probability. Hence $\Pr(\text{no rejection},\,C(D)\ne0)\to0$. Since $c(D)\ne0$ exactly when $C(D)\ne0$, this gives the bound.

It remains to place the two long-run variance rules in the class. Let $u_t=(X_t-\bar X)(Y_t-\bar Y)$, let $\tilde u_t=u_t-\bar u$, and let $\hat\gamma(j)=n^{-1}\sum_{t}\tilde u_{t+j}\tilde u_t$. By Cauchy--Schwarz, $|\hat\gamma(j)|\le\hat\gamma(0)$. Expanding $u_t^2$ as a polynomial in $\bar X$, $\bar Y$, $\xi_t$, $\zeta_t$, $s_t$, $q_t$ and applying Lemma~\ref{lem:ergodic} shows that $\hat\gamma(0)$ converges almost surely to a finite limit.

\emph{Newey--West.} The Bartlett estimate is $\hat v_n=\hat\gamma(0)+2\sum_{1\le j<S_n}(1-j/S_n)\hat\gamma(j)$. It is nonnegative, because the Bartlett weights form a positive-definite sequence. Its weights lie in $[0,1]$, so $\hat v_n\le(2S_n+1)\hat\gamma(0)$. Thus $\hat v_n/n\to0$ whenever $S_n=o_p(n)$. The implemented bandwidth is the \citet{andrews1991hac} AR(1) plug-in $S_n=1.1447(\hat\alpha_n n)^{1/3}$, with $\hat\alpha_n=4\hat\rho_n^2/\{(1-\hat\rho_n)^2(1+\hat\rho_n)^2\}$ and $\hat\rho_n=\sum_t\tilde u_t\tilde u_{t-1}/\sum_t\tilde u_{t-1}^2$. By Lemma~\ref{lem:ergodic} with $L=1$, $\hat\rho_n\to\rho^*$ almost surely. Suppose $|\rho^*|=1$. Then $n^{-1}\sum_t(\tilde u_t\mp\tilde u_{t-1})^2\to0$. By Lemma~\ref{lem:ergodic} this limit equals $P^{-1}\sum_j\E F_j^2$, where $F_j=(\xi_1+s_j)(\zeta_1+q_j)\mp(\xi_0+s_{j-1})(\zeta_0+q_{j-1})+\kappa$ for a constant $\kappa$ that comes from centring. The vector $(\xi_0,\xi_1,\zeta_0,\zeta_1)$ is nondegenerate Gaussian, because $|\gamma(1)|=|2^{2H-1}-1|<1$. $F_j$ is a nonconstant polynomial in it, since it contains $\xi_1\zeta_1$ with coefficient one. So $\E F_j^2>0$, a contradiction. Hence $|\rho^*|<1$, $\hat\alpha_n$ converges to a finite limit, and $S_n=O(n^{1/3})$ almost surely.

\emph{Equal-weighted cosine.} The estimate is $\hat v_n=\nu_n^{-1}\sum_{j=1}^{\nu_n}\Lambda_j^2$ with $\Lambda_j=(2/n)^{1/2}\sum_t\cos\{\pi j(t-1/2)/n\}\tilde u_t$. These cosine vectors, together with $n^{-1/2}\one$, form the orthonormal DCT-II basis of $\mathbb R^n$. So $\sum_{j\le\nu_n}\Lambda_j^2\le\|\tilde u\|^2=n\hat\gamma(0)$ and $\hat v_n/n\le\hat\gamma(0)/\nu_n\to0$ when $\nu_n\to\infty$. The implementation uses $\nu_n=\lfloor0.4n^{2/3}\rfloor$. Its Student-$t$ critical value is bounded by the $\nu=1$ value. \qed

\subsection{Proof of Proposition~\ref{prop:seasm}(iii)}
Fix $(u,v)$ as before. Let $\mu_n=n^{-1}\sum_t\delta_{(X_t,Y_t)}$ be the empirical measure, and define the Gaussian mixture
\[
\mu_{u,v}=\frac1P\sum_{j=0}^{P-1}N\big((s_j,q_j)^\top,\,I_2\big).
\]

\emph{Step 1: $\mu_n\to\mu_{u,v}$ weakly, almost surely.} For fixed $\tau\in\mathbb R^2$, Lemma~\ref{lem:ergodic} with $f_j(w)=\exp\{i\tau^\top(w+(s_j,q_j))\}$ gives convergence of the characteristic function of $\mu_n$ at $\tau$. Intersecting over rational $\tau$ gives one almost-sure event. On it, $n^{-1}\sum_t(X_t^2+Y_t^2)$ also converges, so $(\mu_n)$ is tight. Every subsequential weak limit then has the characteristic function of $\mu_{u,v}$ on a dense set, hence everywhere by continuity. So every subsequential limit equals $\mu_{u,v}$.

\emph{Step 2: the bandwidths converge.} The $X$-marginal of $\mu_{u,v}$ is $\mu^X=P^{-1}\sum_jN(s_j,1)$. As a multiset, $\{s_j\}$ does not depend on $u$, and it equals $\{q_j\}$. Let $G$ be the law of $|X-X'|$ for $X,X'$ independent with law $\mu^X$. It has a continuous, strictly increasing distribution function on $[0,\infty)$, so its median $\sigma^*>0$ is unique. The empirical law of the off-diagonal distances $|X_s-X_t|$, $s\ne t$, is the image of $\mu_n^X\otimes\mu_n^X$ under $(x,x')\mapsto|x-x'|$, minus a diagonal mass $1/n$, renormalized. By Step 1 and the continuous mapping theorem, it converges weakly to $G$. Because $G$ is continuous, the convergence is uniform in the argument, so the median-distance bandwidth $\hat\sigma_n\to\sigma^*$. The positive-distance fallback is eventually inactive. The same holds for $\hat\tau_n$, with the same limit $\sigma^*$.

\emph{Step 3: $V_n\to h(u-v)$.} For bandwidths $(\sigma,\tau)$ and a probability $\nu$ on $\mathbb R^2$, write $k_\sigma(x,x')=\exp\{-(x-x')^2/(2\sigma^2)\}$ and
\[
\HSIC_{\sigma,\tau}(\nu)=\E\,k_\sigma(X_1,X_2)l_\tau(Y_1,Y_2)-2\,\E\,k_\sigma(X_1,X_2)l_\tau(Y_1,Y_3)+\E\,k_\sigma(X_1,X_2)\,\E\,l_\tau(Y_3,Y_4),
\]
with $(X_i,Y_i)$ independent draws from $\nu$ and $l_\tau=k_\tau$. The statistic of Section~\ref{sec:methods} satisfies $V_n=n^{-2}\operatorname{tr}(K^cL^c)=\HSIC_{\hat\sigma_n,\hat\tau_n}(\mu_n)$, the standard identity for the biased estimator \citep{gretton2005measuring}\vc. Each term is an integral of a bounded continuous function against $\nu^{\otimes r}$ with $r\le4$. Weak convergence of $\mu_n$ implies weak convergence of $\mu_n^{\otimes r}$ \citep[Theorem~2.8]{billingsley1999convergence}\vc. So $\HSIC_{\sigma^*,\sigma^*}(\mu_n)\to\HSIC_{\sigma^*,\sigma^*}(\mu_{u,v})$. Next, $|\partial_\sigma k_\sigma|\le(2/e)/\sigma$ uniformly in the arguments. Kernels are bounded by one, so $|\HSIC_{\sigma,\tau}(\nu)-\HSIC_{\sigma',\tau'}(\nu)|\le4(\|k_\sigma-k_{\sigma'}\|_\infty+\|k_\tau-k_{\tau'}\|_\infty)$ for every $\nu$. With Step 2 this gives $V_n\to h(u-v):=\HSIC_{\sigma^*,\sigma^*}(\mu_{u,v})$. The value depends on $(u,v)$ only through $u-v$, by a shift of the index $j$.

\emph{Step 4: $h(d)>0$ for every $d$ when $P\ge5$.} The product $k_\sigma\otimes k_\tau$ is a Gaussian kernel on $\mathbb R^2$. It is characteristic \citep{sriperumbudur2010hilbert}\vc. $\HSIC_{\sigma,\tau}(\nu)$ is the squared MMD between $\nu$ and $\nu^X\otimes\nu^Y$ for this kernel \citep{gretton2005measuring}\vc. So $h(d)=0$ only if $\mu_{u,v}$ equals the product of its marginals. Let $\pi$ be the law of $(s_J,q_J)$ with $J$ uniform on $\{0,\ldots,P-1\}$. Then $\mu_{u,v}=\pi*N(0,I_2)$ and $\mu^X\otimes\mu^Y=(\pi^X\otimes\pi^Y)*N(0,I_2)$. The Gaussian characteristic function has no zeros, so the two are equal only if $\pi=\pi^X\otimes\pi^Y$. Suppose this holds. Let $x^*=\max_js_j$. The equation $a\sin\theta=x^*$ has at most two solutions $\theta\in[0,2\pi)$, so at most two indices $j$ have $s_j=x^*$. Under $\pi$, the point $x^*$ is paired with at most two values of $q$. Under a product measure it is paired with every point of the support of $\pi^Y$. So that support has at most two points. But $q_j$ takes at least $\lceil P/2\rceil\ge3$ distinct values, since each value is taken at most twice. This is a contradiction. Hence $h(d)>0$ for every $d$, and $\min_dh(d)>0$ since $d$ ranges over a finite set.

\emph{Step 5: the rules.} On every path, $T_n/n=V_n\to h(D)>0$ almost surely. If $\kappa_n/n\to0$ in probability, then $\Pr(T_n\le\kappa_n)\le\Pr\{V_n\le h(D)/2\}+\Pr\{\kappa_n/n\geq h(D)/2\}\to0$, using $h(D)\ge\min_dh(d)>0$.

Double-centred Gaussian Gram entries satisfy $|K^c_{ij}|\le2$. So $M=K^c\circ L^c$ has entries bounded by $4$, and it is positive semidefinite by the Schur product theorem.

\emph{Subsampling.} Each window statistic is $T_\ell=\ell^{-1}\sum_{i,j\le\ell}K^c_{ij}L^c_{ij}\le4\ell$. This holds for any bandwidth, including the full-sample bandwidth that the rule reuses. Its $95\%$ quantile is at most $4\ell_n=o(n)$.

\emph{Wild multipliers.} Let $w$ be block-Markov Rademacher with $\Cov(w_s,w_t)=R_{st}=e^{-|s-t|/\ell_n}$, and $\tilde w=Cw$ with $C=I-n^{-1}\one\one^\top$, or $\tilde w=w$ for the raw rule. Given the data, $T^*=n^{-1}\tilde w^\top M\tilde w\ge0$, and
\[
\E_wT^*=n^{-1}\operatorname{tr}(CMC\,R)\le n^{-1}\lambda_{\max}(R)\operatorname{tr}(CMC)\le n^{-1}(1+2\ell_n)\operatorname{tr}(M)\le4(1+2\ell_n).
\]
Here $\lambda_{\max}(R)\le\max_s\sum_tR_{st}\le1+2\ell_n$, $\operatorname{tr}(CMC)\le\operatorname{tr}(M)$ because $\one^\top M\one\ge0$, and $M_{ii}\le4$. With $B$ draws, the rule rejects when $(1+\#\{b:T^*_b\ge T_n\})/(B+1)\le\alpha$, that is, when at most $k_B=\lfloor\alpha(B+1)\rfloor-1$ draws satisfy $T^*_b\ge T_n$. We assume $\alpha(B+1)\ge1$, so $k_B\ge0$. For $B=500$ and $\alpha=0.05$, $k_B=24$. By Markov's inequality each draw exceeds $T_n$ with conditional probability $p_n\le4(1+2\ell_n)/T_n$. So the conditional probability of no rejection is at most $Bp_n$. Since $T_n/\ell_n\to\infty$ almost surely when $\ell_n=o(n)$, dominated convergence gives $\Pr(\text{reject})\to1$. The implementation uses $\ell_n=\lfloor n^{1/2}\rfloor$ for both rules. \qed

\begin{remark}[Circular shifts, proof sketch]\label{rem:shift}
The shift rule compares $T_n$ with $T_n(s)=n\HSIC(\mu_n^{(s)})$, where $\mu_n^{(s)}$ is the empirical measure of $(X_t,Y_{t+s\bmod n})$, for $s$ in a range $S_n$ of about $0.8n$ offsets. When $P$ divides $n$, as for $n=2000$ and $P=50$, the shifted seasonal component of $Y$ is exactly a sinusoid with phase $V+s$. The shifted pair is then a SEAS-M pair with gap $D-s$, and its statistic is also of order $n$. Suppose that $T_n(s)/n\to h(D-s)$ uniformly over $s\in S_n$. Then the leading term of the $p$-value is the fraction of offsets with $h(D-s)\ge h(D)$. That fraction tends to $\pi(D)=P^{-1}\#\{d:h(d)\ge h(D)\}$, because $S_n$ covers each residue class modulo $P$ about equally often. Since $D$ is uniform, $\Pr\{\pi(D)\le\alpha\}\le\alpha$. So the phase gap does not drive the size to one, as it does in part (iii). Two steps are missing. First, uniform convergence over $S_n$ is not proved. Each offset also re-pairs the noise, and the shifted noise has one wrap-around break. Second, $h$ has ties. For even $P$, $h(d)=h(P/2-d)$, because the two limit measures differ by the reflection $y\mapsto-y$. Among offsets in a tied class, the decision is settled by lower-order terms, where the long-memory noise matters. The sketch therefore explains why the size is near $\alpha$ rather than near one. It does not prove that the size equals $\alpha$.
\end{remark}

\begin{remark}[Distance covariance]
The same steps apply to distance covariance with subsampling, with two changes. The distance kernel is unbounded, so Step 3 needs uniform integrability of $|X_s-X_t|$ under $\mu_n^{\otimes2}$. This follows from almost-sure convergence of $n^{-1}\sum_tX_t^2$. The window bound $T_\ell\le4\ell$ becomes a bound by $\ell$ times the largest window second moment. That moment grows at most logarithmically for Gaussian noise. We have not written out these two steps, so this case is stated without proof.
\end{remark}

\subsection{Proof of Proposition~\ref{prop:halfmean}}
\emph{(i)} Let $S_k=\sum_{t\le k}\xi_t$. Unit-variance fGn has $\Var(S_k)=k^{2H}$ exactly. Write $A=S_m$ and $B=S_{2m}-S_m$, so $\Var(A)=\Var(B)=m^{2H}$ by stationarity. From $\Var(A+B)=(2m)^{2H}$ we get $2\Cov(A,B)=(2m)^{2H}-2m^{2H}$. Hence
\[
\Var(A-B)=2m^{2H}-2\Cov(A,B)=4m^{2H}-(2m)^{2H},
\]
and $\Delta_m=(A-B)/m$ gives \eqref{eq:halfvar}. The ratio $R_m$ is $\delta/\Var(\Delta_m)^{1/2}$.

\emph{(ii)} The first-half mean minus the second-half mean of $Z$ equals $\Delta_m-\delta$. It is Gaussian with mean $-\delta$ and variance $\Var(\Delta_m)$, and $\mu$ cancels. The power formula is the two-sided normal tail at noncentrality $R_m$.

\emph{(iii)} Let $L$ be any $(n-1)\times n$ matrix of rank $n-1$ with $L\one=0$. The statistic $LZ$ is a maximal invariant under $Z\mapsto Z+b\one$. So every invariant test is a function of $LZ$ \citep[Theorem~6.2.1]{lehmann2005testing}\vc. Its law is $N(\delta Lc,\Sigma)$ with $\Sigma=L\Gamma L^\top$, which does not involve $\mu$. By the Neyman--Pearson lemma, the most powerful level-$\alpha$ test of $\delta=0$ against a fixed $\delta>0$ in this family rejects for large $(Lc)^\top\Sigma^{-1}LZ$. Its power is $1-\Phi\{z_{1-\alpha}-\delta(c^\top L^\top\Sigma^{-1}Lc)^{1/2}\}$. Any invariant level-$\alpha$ test has at most this power.

It remains to show $L^\top(L\Gamma L^\top)^{-1}L=\Gamma^{-1}-\Gamma^{-1}\one(\one^\top\Gamma^{-1}\one)^{-1}\one^\top\Gamma^{-1}$. Multiply the left side by $\Gamma^{1/2}$ on both sides. The result is the orthogonal projection onto the range of $\Gamma^{1/2}L^\top$. That range has dimension $n-1$ and is orthogonal to $\Gamma^{-1/2}\one$, since $L\Gamma^{1/2}\Gamma^{-1/2}\one=L\one=0$. So the projection equals $I-\Gamma^{-1/2}\one\one^\top\Gamma^{-1/2}/(\one^\top\Gamma^{-1}\one)$. Multiplying back by $\Gamma^{-1/2}$ on both sides gives the identity, and hence $c^\top L^\top\Sigma^{-1}Lc=I_n$. By the Gauss--Markov theorem for the regression of $Z$ on $(\one,c)$ with covariance $\Gamma$, $I_n^{-1}$ is the variance of the generalized-least-squares estimator of $\delta$. \qed

Since that estimator is the best linear unbiased one, $I_n\ge1/\Var(\Delta_m)$, and the cap in (iii) is never below the power in (ii). Part (iii) assumes that $H$ and the unit variance are known. A test that controls size over a range of $H$ is also level $\alpha$ at each $H$ in the range, so the cap applies to it too.
